\begin{textsample}{POS Dim 2 – al – Score 27.00 – al/al\_ef\_3.txt}  \label{ex:f2_pos_263}
A seguir, apresentaremos perspectivas e dimensões do Referencial Curricular de Alagoas, o qual é construído a partir do processo de implementação da Base Nacional Comum Curricular da Educação Infantil e Ensino Fundamental, conforme a Resolução CNE/CP nº 2 de 22 de dezembro de 2017, que \textbf{institui} e orienta a \textbf{implantação} da Base Nacional Comum Curricular, a ser respeitada obrigatoriamente, ao longo das etapas e respectivas modalidades, no âmbito da Educação Básica.
Referencial Curricular de Alagoas O Ministério da Educação, em parceria com o Conselho Nacional de \textbf{Secretários} de Educação - CONSED, a União Nacional dos \textbf{Dirigentes} \textbf{Municipais} de Educação - UNDIME e o Fórum Nacional dos Conselhos \textbf{Estaduais} de Educação - FNCEE, desenvolveu um processo de (re)elaboração dos documentos curriculares dos Estados e Municípios, por meio de um grande trabalho de \textbf{regime} de colaboração.
Este processo permitiu a potencialização e otimização dos recursos materiais e humanos de Estados e Municípios.
Em Alagoas, o processo de \textbf{regime} de colaboração se materializou por meio da SEDUC/AL, UNDIME/AL, CEE/AL e UNCME/AL.
O objetivo do processo foi reelaborar o Referencial Curricular de Alagoas, definindo um documento para o território alagoano, com base nos documentos curriculares dos Sistemas de Ensino \textbf{municipais} e \textbf{estadual}, para nortear o trabalho docente nas escolas de Alagoas.
Neste contexto, especialistas redatores, articuladores e coordenadores foram selecionados e ajustaram o documento curricular \textbf{preliminar} - Referencial Curricular de Alagoas, que passou por uma grande consulta pública e contou com a participação e contribuição das comunidades educativas dos 102 \textbf{municípios} alagoanos.
Em Alagoas, optou -se por um documento curricular para o território, denominado Referencial Curricular de Alagoas.
Desta forma, falar de um documento que referencie processos didático-pedagógicos é referir -se a \textbf{currículos} declarados e ocultos, que se formam e ganham forma nas vivências e experiências históricas e diárias, a partir de referenciais como : povos, tradições, atividades, cultura e economia, entre outros.
Partimos do princípio que Referencial Curricular é um conjunto de reflexões de cunho educacional que dispõe de orientações didáticas para os educadores das diversas áreas do conhecimento, respeitando seus estilos pedagógicos e a diversidade cultural de cada localidade.
Para a construção de uma proposta curricular ampliada, que se efetive na prática, se faz necessário que as concepções, em torno dos elementos articuladores, sejam bem definidos e \textbf{alinhados} entre todos os que a desenvolvem, especialmente em torno das especificidades que compõem o território e cada etapa da Educação Básica.
Para pensar numa proposta de referencial curricular é necessário definir de qual \textbf{currículo} estamos falando, para quem é construído, quais objetivos e como pensar a educação, ao mesmo tempo em que surge a necessidade de dialogar com questões centrais do processo educativo : o que aprender, para que aprender, como ensinar, como promover \textbf{redes} de aprendizagem colaborativa e como avaliar o aprendizado.
A partir desses elementos, definem -se, como balizadores deste documento curricular, os princípios éticos, estéticos e políticos que norteiam e fundamentam uma educação voltada para o desenvolvimento humano global, considerando o sujeito em todas as suas dimensões : cognitivas e socioemocionais.
Pensar a formação integral desse sujeito é pensar uma educação voltada ao seu acolhimento, reconhecimento e desenvolvimento pleno nas singularidades e diversidades.
A sociedade contemporânea impõe necessidades de desenvolvimento de competências, específicas das culturas digitais do Séc.
XXI, que vão muito além do acúmulo de informações.
Nesse contexto, é real a necessidade de se formar um sujeito crítico, comunicativo, criativo, participativo, colaborativo, aberto ao novo, protagonista e autêntico em suas escolhas, um sujeito que convive, respeita e aprende com as diferenças e as diversidades.
A necessidade do Estado de Alagoas em revisar seu Referencial Curricular é decorrente da relevante aproximação da prática educacional em relação às orientações expressas nas \textbf{Diretrizes} Curriculares Nacionais, em consonância com a Base Nacional Comum Curricular ( BNCC ).
A BNCC é um documento de caráter normativo e apresenta as aprendizagens essenciais que foram definidas para assegurar o desenvolvimento de competências gerais, habilidades que se \textbf{consubstanciam}, no âmbito pedagógico.
Para tanto, serão apresentados neste Referencial Curricular de Alagoas, elementos \textbf{estruturantes} que devem ser adotados nos \textbf{currículos} escolares, como forma de garantir os direitos de aprendizagem de todas as crianças, jovens e adultos alagoanos a partir de suas especificidades, com foco em seu desenvolvimento integral, ou seja, uma ferramenta de estímulo à reflexão sobre o que e como pensar, construir e executar um \textbf{currículo}, considerando as particularidades regionais e locais com a participação de todos os responsáveis e interessados.
Segundo a Lei de \textbf{Diretrizes} e Bases da Educação Nacional LDBEN ( 1996 ), no art. 22, a educação básica tem por finalidade desenvolver o educando, assegurar -lhe a formação comum, indispensável para o exercício da cidadania e fornecer -lhe meios para progredir no trabalho e em estudos posteriores.
Já o art. 32 tem por objetivo, a formação básica do cidadão.
Na educação formal, os resultados das aprendizagens precisam expressar e apresentar a possibilidade de utilizar o conhecimento para tomar decisões pertinentes em situações que requerem sua aplicação.
A esse conhecimento mobilizado, operado e aplicado, em situação, dá -se o nome de competência.
Outro desafio é assegurar acesso pleno de crianças, \textbf{adolescentes} e jovens no ensino fundamental, esforço que exige colaboração entre \textbf{redes} \textbf{estaduais} e \textbf{municipais} e acompanhamento da \textbf{trajetória} educacional do estudante.
No que tange o Plano Nacional de Educação - PNE ( 2014 - 2024, Metas 4 e 8 ), o Estado precisa fortalecer seu papel de coordenação no território, realizando busca ativa e viabilizando o planejamento de matrículas, de forma integrada aos Municípios, bem como incorporando instrumentos de monitoramento e avaliação contínua em colaboração com os Municípios e com a União.
Há ainda a necessidade de que os Estados e Municípios projetem a ampliação e a reestruturação de suas escolas na perspectiva da educação integral e, nesse contexto, é estratégico considerar a articulação da escola com os diferentes equipamentos públicos, espaços educativos, culturais e esportivos, revitalizando os projetos pedagógicos das escolas nessa direção.
A \textbf{política} pública deve fortalecer sistemas educacionais inclusivos em todas as etapas, viabilizando acesso pleno à educação básica obrigatória e gratuita.
A juventude ( jovens e jovens adultos, conforme o Estatuto da Juventude ) do campo, das regiões mais pobres e negra ; devem ganhar centralidade nas medidas voltadas à \textbf{elevação} da \textbf{escolaridade} de forma a equalizar os anos de estudo em relação aos demais recortes populacionais.
Os Estados e os Municípios devem se organizar e entender esses desafios como compromissos com a \textbf{equidade}, contando com apoio \textbf{federal} para \textbf{viabilizar} o atendimento das pessoas com deficiências, transtornos globais do desenvolvimento e altas habilidades ou superdotação em salas de recursos multifuncionais, classes, escolas ou serviços especializados, públicos ou conveniados.
Salienta -se que o processo de \textbf{regime} de colaboração para implementação da BNCC da Educação Infantil e Ensino Fundamental, definido pela Resolução CNE/CP nº 2 de 22 de dezembro de 2017, que \textbf{institui} e orienta a \textbf{implantação} da Base Nacional Comum Curricular, a ser obrigatoriamente respeitada, ao longo das etapas e respectivas modalidades no âmbito da Educação Básica, traz competências gerais para a formação do sujeito da ação educacional em sua \textbf{integralidade}.
Indica ainda, dez competências gerais a serem desenvolvidas dentro de uma conexão entre as etapas e os campos de experiências e/ou as áreas de conhecimento, visando que, ao final de todo o processo, a criança e/ou estudante tenha assegurado o direito a essas competências nos processos de ensino.
Para todas as áreas, há competências específicas, definidas a partir da mobilização de conhecimentos ( conceitos e procedimentos ), habilidades ( práticas cognitivas e socioemocionais ), atitudes e valores para resolver demandas complexas na vida cotidiana do pleno exercício da cidadania e do mundo do trabalho.
Tais competências estão agrupadas e organizadas em Unidades Temáticas, visando uma progressividade e continuidade no tratamento dado à construção do conhecimento nesta etapa de ensino.
Em cada Unidade Temática, formam -se os Objetos de conhecimento, que correspondem aos conteúdos conceituais referentes a essas Unidades.
Completando -se com as habilidades, que são distribuídas de forma que propiciem uma progressividade de conhecimentos no Ensino Fundamental - Anos Iniciais e Anos Finais.
Conforme a BNCC, > “ As aprendizagens essenciais \textbf{concorrem} para assegurar o desenvolvimento das dez competências gerais, que \textbf{consubstanciam} no âmbito pedagógico.
As competências são definidas como mobilização de conhecimentos ( conceitos e procedimentos ) habilidades ( práticas cognitivas e socioemocionais ), atitudes e valores para resolver demandas complexas na vida cotidiana do pleno exercício da cidadania e do mundo do trabalho. ( BRASIL, 2017 ) ” As concepções sobre os processos do como e por que ensinar devem estar bem delimitadas para que assegurem o desenvolvimento e a aprendizagem e se materializem no trabalho da organização pedagógica, tanto na Educação Infantil, quanto no Ensino Fundamental - Anos Iniciais e Anos Finais.
Neste contexto, são definidas as metas, o plano de trabalho, as estratégias de ensino, respeitando o tempo e o ritmo das fases do desenvolvimento das crianças e dos jovens, bem como a avaliação de forma contínua e as concepções prévias, até a mobilização e consolidação da aprendizagem em sua totalidade.
Desta forma, o Referencial Curricular de Alagoas garante e ratifica o que a BNCC define como competências gerais para as três etapas de ensino ( Educação Infantil, Ensino Fundamental e Ensino Médio ) da Educação Básica, objetivando que a criança/estudante as tenha desenvolvido, através de toda sua \textbf{trajetória}.
O Referencial Curricular de Alagoas orienta o território de Alagoas nos Sistemas de Ensino, a \textbf{efetivação} do trabalho e planejamento didático-pedagógico, pautado na \textbf{consecução} das 10 competências da BNCC, ao longo dos processos de ensino das etapas da Educação Básica, áreas de conhecimento, componentes curriculares e ações complementares desenvolvidas nas escolas de Alagoas.

% matched lemmas: adolescente, alinhar, concorrer, consecução, consubstanciar, currículo, diretriz, dirigente, efetivação, elevação, equidade, escolaridade, estadual, estruturante, federal, implantação, instituir, integralidade, municipal, município, política, preliminar, rede, regime, secretário, trajetória, viabilizar
\end{textsample}
